% =====================================================================
\section{Wi-Fi и ESPUI: управляем платой с телефона}\label{les:radio}
% =====================================================================
\lessonintro
  {подключить ESP32-S3 к Wi-Fi и собрать панель управления светодиодом, которая открывается на телефоне}
  {урок 1 (светодиод), урок 4 (ШИМ), урок 2 (Serial Monitor)}
  {светодиодом можно управлять с телефона: переключатель вкл/выкл и ползунок яркости}
  {Wi-Fi-сеть на \SI{2.4}{GHz}; светодиод с резистором на GPIO4}


\subsection{Что такое радиоволны}

Радиоволны --- это \textbf{электромагнитные волны}, как и свет. Отличаются они только
\emph{частотой}: сколько раз в секунду колеблется волна. Мы уже знаем, что частоту измеряют в герцах.
Для радио получаются огромные числа, поэтому используют приставки:
1~килогерц (кГц) = 1000~Гц, 1~мегагерц (МГц) = миллион герц, 1~гигагерц (ГГц) = миллиард герц.

\begin{figure}[H]
\centering
\begin{tikzpicture}[x=0.62cm,y=1cm]
  % шкала частот: логарифм, 10^4 .. 10^20 Гц -> x = 0..16 (log10 f - 4)
  \shade[left color=blkrtc!60,right color=blkrf] (0,0) rectangle (8.5,1);
  \shade[left color=blkrf,right color=esp!25] (8.5,0) rectangle (10.4,1);
  \fill[esp!25] (10.4,0) rectangle (10.55,1);
  \shade[left color=red,right color=yellow] (10.55,0) rectangle (10.63,1);
  \shade[left color=yellow,right color=blue!80] (10.63,0) rectangle (10.72,1);
  \shade[left color=blue!80,right color=violet] (10.72,0) rectangle (10.8,1);
  \shade[left color=violet!40,right color=gray!30] (10.8,0) rectangle (16,1);
  \draw (0,0) rectangle (16,1);
  % волна с убывающей длиной
  \draw[thick,accent,domain=0:16,samples=500,smooth] plot (\x,{1.6+0.28*sin(0.35*\x*\x r)});
  \foreach \x/\l in {0/10\,кГц,2/1\,МГц,5/1\,ГГц,8/1\,ТГц,11/$10^{15}$,14/$10^{18}$}
    \draw (\x,0) -- (\x,-0.12) node[below,font=\scriptsize]{\l};
  \node[font=\sffamily\scriptsize,align=center] at (2.2,0.5) {радио};
  \node[font=\sffamily\scriptsize,align=center] at (6.6,0.5) {СВЧ};
  \node[font=\sffamily\scriptsize,align=center] at (9.4,0.5) {ИК};
  \node[font=\sffamily\scriptsize,align=center] at (12.3,0.5) {УФ};
  \node[font=\sffamily\scriptsize,align=center] at (14.6,0.5) {рентген};
  % примеры: номера на шкале и расшифровка ниже
  \tikzset{mk/.style={circle,draw=black,fill=white,inner sep=0.8pt,font=\sffamily\bfseries\tiny,minimum size=3.6mm}}
  \foreach \x/\n/\yy in {1.1/1/-0.75,4.0/2/-0.75,4.95/3/-1.35,5.38/4/-0.75,5.7/5/-1.95,9.5/6/-0.75,10.65/7/-1.35,14.3/8/-0.75}{
    \draw[thin,gray] (\x,0) -- (\x,\yy);
    \node[mk] at (\x,\yy-0.18) {\n};
  }
  \node[mk,fill=esp!25] at (5.38,-0.93) {4};
  \node[anchor=north west,font=\sffamily\scriptsize,align=left] at (0,-2.6) {%
    \textbf{1} АМ-радио, 0,5--1,6 МГц\qquad \textbf{2} FM-радио, 100 МГц\qquad \textbf{3} мобильная связь, 0,9--2 ГГц\\[1pt]
    \textbf{\textcolor{esp}{4} Wi-Fi и Bluetooth, 2,4 ГГц} (и микроволновка)\qquad \textbf{5} Wi-Fi 5 ГГц\\[1pt]
    \textbf{6} пульт от телевизора (инфракрасный)\qquad \textbf{7} видимый свет\qquad \textbf{8} рентгеновский снимок};
  \draw[{Stealth[length=2mm]}-{Stealth[length=2mm]}] (0,2.25) -- node[above,font=\scriptsize]{частота растёт, длина волны уменьшается} (16,2.25);
\end{tikzpicture}
\caption{Шкала электромагнитных волн (частоты отложены так, что каждое деление больше предыдущего
в тысячу раз). Видимый свет занимает крошечную полоску, а Wi-Fi --- соседи микроволновки}
\end{figure}

Все электромагнитные волны летят со скоростью света $c = 300\,000$~км/с. Зная частоту, можно найти
\textbf{длину волны} --- расстояние между соседними «гребнями»:
\[
  \lambda = \frac{c}{f} = \frac{300\,000\,000~\text{м/с}}{2\,400\,000\,000~\text{Гц}} = 0{,}125~\text{м} = 12{,}5~\text{см}.
\]
Антенна обычно делается длиной в четверть волны --- около 3~см. Посмотри на зигзаг на модуле ESP32-S3:
если его «распрямить», получится примерно такая длина!

\begin{fact}
Микроволновая печь греет еду волнами почти той же частоты, что и Wi-Fi, --- 2,45~ГГц. Только
мощность у неё около 800~Вт, а у Wi-Fi платы --- меньше 0,1~Вт, в десять тысяч раз слабее. Поэтому
работающая микроволновка может мешать Wi-Fi, но Wi-Fi никого не нагреет.
\end{fact}

\subsection{Wi-Fi и его каналы}

Wi-Fi --- это способ передавать цифровые данные по радио. Биты «прячутся» в радиоволне: передатчик
очень быстро меняет её параметры, а приёмник расшифровывает изменения обратно в нули и единицы.

Диапазон 2,4~ГГц поделён на \textbf{каналы}, как дорога --- на полосы. Если две сети работают на одном
канале, они мешают друг другу. Каналы шириной 20~МГц частично перекрываются, поэтому не мешают друг
другу только каналы 1, 6 и 11:

\begin{figure}[H]
\centering
\begin{tikzpicture}[x=0.19cm,y=1cm]
  \draw[->] (-2,0) -- (78,0) node[right,font=\scriptsize]{МГц};
  \foreach \ch in {1,...,13}{
    \pgfmathsetmacro{\c}{2412+5*(\ch-1)-2400}
    \pgfmathsetmacro{\hh}{(\ch==1||\ch==6||\ch==11)?1.5:1.0}
    \ifnum\ch=1 \def\col{esp}\else\ifnum\ch=6 \def\col{okgreen}\else\ifnum\ch=11 \def\col{accent}\else\def\col{gray!45}\fi\fi\fi
    \draw[thick,\col] (\c-10,0) .. controls (\c-7,\hh) and (\c+7,\hh) .. (\c+10,0);
    \node[font=\scriptsize,text=\col] at (\c,-0.3) {\ch};
  }
  \foreach \f in {2400,2420,2440,2460,2480}{\pgfmathsetmacro{\x}{\f-2400}\draw (\x,0.05) -- (\x,-0.05) node[below=3mm,font=\scriptsize]{\f};}
  \node[font=\sffamily\scriptsize,anchor=west] at (-2,1.7) {номера каналов --- под осью; не перекрываются 1, 6 и 11};
\end{tikzpicture}
\caption{Каналы Wi-Fi в диапазоне 2,4~ГГц}
\end{figure}

Начнём с того, что посмотрим, какие сети «видит» плата. Эта программа выводит имя каждой сети,
её канал и силу сигнала:

\codecap{Поиск Wi-Fi сетей}
\codeinput{l07_scan}

Сила сигнала (RSSI) измеряется в \textbf{дБм} (децибел-милливаттах) и всегда отрицательна: чем
ближе к нулю, тем сильнее сигнал.

\begin{figure}[H]
\centering
\begin{tikzpicture}[x=0.16cm]
  \shade[left color=okgreen!70,right color=esp!70] (0,0) rectangle (70,0.6);
  \foreach \v in {-30,-40,...,-100}{\pgfmathsetmacro{\x}{-(\v)-30}\draw (\x,0) -- (\x,-0.12) node[below,font=\scriptsize]{\v};}
  \foreach \x/\t in {5/отлично: рядом с роутером,22/хорошо,37/средне: страницы\\ открываются медленно,55/плохо: связь рвётся}
    \node[font=\sffamily\scriptsize,align=center,above] at (\x,0.65) {\t};
  \node[font=\sffamily\scriptsize,anchor=west] at (71,0.3) {дБм};
\end{tikzpicture}
\caption{Как понимать силу сигнала RSSI}
\end{figure}

\subsection{Bluetooth}

\textbf{Bluetooth} работает в том же диапазоне 2,4~ГГц, но создан для другого: связывать
\emph{близкие} устройства (наушники, часы, телефон) и при этом тратить очень мало энергии.
ESP32-S3 поддерживает \textbf{Bluetooth Low Energy} (BLE, «Bluetooth с низким энергопотреблением»).

\begin{table}[H]
\centering\small
\begin{tabularx}{\textwidth}{@{}l X X@{}}
\toprule
 & \textbf{Wi-Fi} & \textbf{Bluetooth LE}\\
\midrule
Дальность & десятки метров & 10--30 метров\\
Скорость & миллионы бит в секунду & сотни тысяч бит в секунду\\
Потребление & большое (до 300~мА в пике) & маленькое (можно годами работать от батарейки)\\
Нужен роутер? & обычно да & нет, устройства связываются напрямую\\
Для чего & интернет, веб-страницы, видео & датчики, фитнес-браслеты, умные замки, наушники\\
\bottomrule
\end{tabularx}
\caption{Wi-Fi и Bluetooth LE: одна частота --- разные задачи}
\end{table}

Устройства BLE постоянно «кричат» в эфир короткие сообщения: «Я тут, меня зовут так-то». Послушаем их:

\codecap{Поиск Bluetooth-устройств поблизости}
\codeinput{l07_ble_scan}

Сканер показывает только устройства, которые передают объявления BLE. Подключение и обмен
данными разобраны в уроке 12.

\subsection{Режимы Wi-Fi}

\begin{figure}[H]
\centering
\begin{tikzpicture}[node distance=10mm]
  \begin{scope}
    \node[blk,fill=blkrf] (r) {Роутер};
    \node[blk,fill=blkcpu,below left=10mm and 0mm of r] (e) {ESP32-S3\\ \scriptsize STA};
    \node[blk,fill=blkper,below right=10mm and 0mm of r] (p) {Телефон/ПК};
    \draw[darr,dashed] (e) -- (r); \draw[darr,dashed] (p) -- (r);
    \node[font=\sffamily\bfseries\small,above=2mm of r] {Станция (STA)};
  \end{scope}
  \begin{scope}[xshift=6cm]
    \node[blk,fill=blkcpu] (e2) {ESP32-S3\\ \scriptsize AP 192.168.4.1};
    \node[blk,fill=blkper,below left=10mm and -2mm of e2] (p2) {Телефон};
    \node[blk,fill=blkper,below right=10mm and -2mm of e2] (p3) {Ноутбук};
    \draw[darr,dashed] (p2) -- (e2); \draw[darr,dashed] (p3) -- (e2);
    \node[font=\sffamily\bfseries\small,above=2mm of e2] {Точка доступа (AP)};
  \end{scope}
  \begin{scope}[xshift=12cm]
    \node[blk,fill=blkrf] (r3) {Роутер};
    \node[blk,fill=blkcpu,below=10mm of r3] (e3) {ESP32-S3\\ \scriptsize AP + STA};
    \node[blk,fill=blkper,below=6mm of e3] (p4) {Телефон};
    \draw[darr,dashed] (e3) -- (r3); \draw[darr,dashed] (p4) -- (e3);
    \node[font=\sffamily\bfseries\small,above=2mm of r3] {Смешанный};
  \end{scope}
\end{tikzpicture}
\caption{Три режима работы Wi-Fi. В опытах используем режим станции}
\end{figure}

Стек Wi-Fi --- это несколько задач FreeRTOS, которые по умолчанию работают на ядре~0. Наши задачи
закреплены за ядром~1, поэтому сеть не мешает кнопке и экрану.

\subsection{Подключение к сети}

\codecap{Подключение к Wi-Fi}
\codeinput{l8_connect}

\begin{caution}
ESP32-S3 работает только в диапазоне \SI{2.4}{GHz}. Сети \SI{5}{GHz} он не увидит.
\end{caution}

\subsubsection{Учимся программировать: цикл \texttt{while} и логическое «и»}

Подключение к сети занимает несколько секунд, и заранее неизвестно, сколько именно. Цикл \code{for} здесь
не подходит --- он для случаев, когда число повторений известно. Нужен цикл \code{while} («пока»): он
повторяет команды, \emph{пока} условие верно.

\begin{figure}[H]
\centering
\begin{tikzpicture}[node distance=7mm]
  \node[fcproc] (b) {WiFi.begin(SSID, PASS)};
  \node[fcdec,below=of b] (q) {ещё не подключились\\ \textbf{и} прошло < 15 с ?};
  \node[fcproc,below=9mm of q] (w) {подождать 300 мс,\\ напечатать «.»};
  \node[fcdec,right=18mm of q] (ok) {подключились?};
  \node[fcproc,above=7mm of ok] (y) {напечатать IP};
  \node[fcproc,below=9mm of ok] (n) {«Не получилось»};
  \draw[arr] (b) -- (q);
  \draw[arr] (q) -- node[left,font=\sffamily\scriptsize]{да} (w);
  \draw[arr] (w.west) -- ++(-8mm,0) |- (q.west);
  \draw[arr] (q) -- node[above,font=\sffamily\scriptsize]{нет} (ok);
  \draw[arr] (ok) -- node[right,font=\sffamily\scriptsize]{да} (y);
  \draw[arr] (ok) -- node[right,font=\sffamily\scriptsize]{нет} (n);
\end{tikzpicture}
\caption{Блок-схема подключения: ждём в цикле, но не дольше 15 секунд}
\end{figure}

В условии цикла два вопроса соединены знаком \code{\&\&} --- логическим «\textbf{и}»: цикл продолжается, только
если верны \emph{оба} условия. Как только плата подключилась \emph{или} время вышло, цикл заканчивается.
Без второго условия плата могла бы ждать вечно, если пароль неверный.

\begin{caution}
Если условие \code{while} никогда не станет ложным, программа застрянет в цикле навсегда. Проверяй, что
рано или поздно цикл закончится. Выйти из любого цикла досрочно можно командой \code{break}, а из
функции --- командой \code{return} (так сделано в программе, если подключиться не удалось).
\end{caution}

\begin{caution}
ESP32-S3 работает только в диапазоне 2,4~ГГц. Сеть 5~ГГц он не увидит. Если роутер раздаёт обе сети с
одинаковым именем, иногда помогает включить на телефоне точку доступа и подключить плату к нему.
\end{caution}


\subsection{MAC-адрес и IP-адрес}

Когда в сети много устройств, данные нужно доставлять точно по адресу. Для этого у каждого
устройства есть \emph{два} адреса.

\begin{figure}[H]
\centering
\begin{tikzpicture}[x=1cm,y=1cm]
  \node[font=\sffamily\bfseries,anchor=west] at (0,1.9) {MAC-адрес --- «номер паспорта»};
  \foreach \b/\c [count=\i] in {7C/blkrtc,DF/blkrtc,A1/blkrtc,0F/blkper,3E/blkper,20/blkper}{
    \node[draw,fill=\c,minimum width=11mm,minimum height=8mm,font=\ttfamily] at (\i*1.25,1) {\b};
    \ifnum\i<6 \node[font=\ttfamily] at (\i*1.25+0.625,1) {:};\fi
  }
  \draw[decorate,decoration={brace,mirror,amplitude=4pt}] (0.7,0.5) -- (3.95,0.5) node[midway,below=5pt,font=\sffamily\scriptsize,align=center]{производитель\\ (здесь --- Espressif)};
  \draw[decorate,decoration={brace,mirror,amplitude=4pt}] (4.45,0.5) -- (7.7,0.5) node[midway,below=5pt,font=\sffamily\scriptsize,align=center]{номер конкретного\\ устройства};
  \node[font=\sffamily\scriptsize,align=left,anchor=west] at (8.4,1) {6 байт, записаны в\\ шестнадцатеричной системе.\\ Даётся на заводе, у всех разный.};
  \begin{scope}[yshift=-3.3cm]
  \node[font=\sffamily\bfseries,anchor=west] at (0,1.9) {IP-адрес --- «почтовый адрес»};
  \foreach \b/\c [count=\i] in {192/blkmem,168/blkmem,1/blkmem,50/blkrf}{
    \node[draw,fill=\c,minimum width=13mm,minimum height=8mm,font=\ttfamily] at (\i*1.6,1) {\b};
    \ifnum\i<4 \node[font=\ttfamily] at (\i*1.6+0.8,1) {.};\fi
  }
  \draw[decorate,decoration={brace,mirror,amplitude=4pt}] (1,0.5) -- (5.4,0.5) node[midway,below=5pt,font=\sffamily\scriptsize,align=center]{номер сети\\ («город и улица»)};
  \draw[decorate,decoration={brace,mirror,amplitude=4pt}] (5.8,0.5) -- (7.0,0.5) node[midway,below=5pt,font=\sffamily\scriptsize,align=center]{номер\\ устройства};
  \node[font=\sffamily\scriptsize,align=left,anchor=west] at (8.4,1) {4 числа от 0 до 255.\\ Выдаёт роутер, когда плата\\ подключается к сети.};
  \end{scope}
\end{tikzpicture}
\caption{Два адреса каждого сетевого устройства}
\end{figure}

\begin{itemize}
  \item \textbf{MAC-адрес} --- постоянный «номер паспорта» сетевого устройства. Его записывают на заводе,
        и он не меняется при переезде в другую сеть. Роутер по MAC-адресу узнаёт устройство
        «в лицо»: например, можно запретить подключение незнакомым устройствам или выдать своему
        всегда один и тот же IP.
  \item \textbf{IP-адрес} --- «почтовый адрес» в конкретной сети. В домашней сети он обычно
        выглядит как \code{192.168.1.\emph{xx}} или \code{192.168.0.\emph{xx}}. Переедешь в другую сеть
        (например, в школьную) --- адрес будет другим, как у человека при переезде.
\end{itemize}

\begin{figure}[H]
\centering
\begin{tikzpicture}[node distance=6mm and 8mm]
  \node[blk,fill=blkrf,minimum width=30mm,minimum height=12mm] (r) {Роутер\\ \scriptsize\ttfamily 192.168.1.1};
  \node[blk,fill=gray!15,right=24mm of r,minimum height=12mm] (net) {Интернет};
  \node[blk,fill=blkcpu,below left=12mm and 4mm of r,minimum width=30mm] (e) {ESP32-S3\\ \scriptsize\ttfamily 192.168.1.50\\ \scriptsize\ttfamily 7C:DF:A1:0F:3E:20};
  \node[blk,fill=blkmem,below=12mm of r,minimum width=30mm] (t) {Телефон\\ \scriptsize\ttfamily 192.168.1.23\\ \scriptsize\ttfamily 5A:11:9B:\ldots};
  \node[blk,fill=blkmem,below right=12mm and 4mm of r,minimum width=30mm] (n) {Ноутбук\\ \scriptsize\ttfamily 192.168.1.37\\ \scriptsize\ttfamily 3C:22:FB:\ldots};
  \draw[darr,dashed] (e) -- (r); \draw[darr,dashed] (t) -- (r); \draw[darr,dashed] (n) -- (r);
  \draw[darr] (r) -- (net);
  \node[font=\sffamily\scriptsize,align=center,above=1mm of net] {для интернета вся\\ квартира --- один адрес};
\end{tikzpicture}
\caption{Домашняя сеть: у каждого устройства свой IP-адрес внутри сети и свой MAC-адрес}
\end{figure}

\subsubsection{Кто раздаёт IP-адреса}

Когда плата подключается, она не знает своего IP. Она громко спрашивает всю сеть, и роутер
выдаёт ей свободный адрес. Этот механизм называется \textbf{DHCP}:

\begin{figure}[H]
\centering
\begin{tikzpicture}[y=-0.8cm]
  \node[blk,fill=blkcpu,minimum width=28mm] (e) at (0,0) {ESP32-S3};
  \node[blk,fill=blkrf,minimum width=28mm] (r) at (9,0) {Роутер};
  \draw[thick,gray] (e.south) -- (0,5.2);
  \draw[thick,gray] (r.south) -- (9,5.2);
  \draw[arr,accent] (0,1.3) -- node[above,font=\sffamily\scriptsize]{«Кто тут раздаёт адреса? Мой MAC --- 7C:DF:\ldots»} (9,1.3);
  \draw[arr,okgreen] (9,2.4) -- node[above,font=\sffamily\scriptsize]{«Я! Могу дать тебе 192.168.1.50»} (0,2.4);
  \draw[arr,accent] (0,3.5) -- node[above,font=\sffamily\scriptsize]{«Беру 192.168.1.50»} (9,3.5);
  \draw[arr,okgreen] (9,4.6) -- node[above,font=\sffamily\scriptsize]{«Договорились. Роутер --- 192.168.1.1»} (0,4.6);
\end{tikzpicture}
\caption{Как плата получает IP-адрес (DHCP)}
\end{figure}

Неудобство в том, что при следующем включении роутер может выдать \emph{другой} адрес. Решений три:
\begin{enumerate}
  \item попросить роутер всегда выдавать этому MAC-адресу один и тот же IP (в настройках роутера);
  \item задать \textbf{постоянный (статический) IP} прямо в программе;
  \item вообще не запоминать IP, а обращаться к плате \textbf{по имени} --- об этом ниже.
\end{enumerate}

\subsubsection{Как задать IP, MAC и имя в программе}

\codecap{Постоянный IP-адрес, свой MAC и имя в сети}
\codeinput{l07_names}

\begin{itemize}
  \item \code{WiFi.config(ip, gateway, subnet, dns)} задаёт IP вручную. Выбирай адрес из своей сети
        (первые три числа --- как у роутера) и такой, который точно никем не занят. Если адрес совпадёт
        с чужим, оба устройства будут работать с ошибками.
  \item \code{esp\_wifi\_set\_mac()} меняет MAC-адрес. Обычно этого делать не нужно; первый байт
        \code{0x02} означает «адрес назначен вручную», это локально назначенный адрес; его уникальность внутри сети нужно обеспечить самостоятельно.
  \item \code{WiFi.setHostname()} --- имя, которое будет видно в списке устройств на странице роутера.
\end{itemize}

\subsection{Имя вместо адреса: mDNS}

Запоминать \code{192.168.1.50} неудобно. В интернете эту задачу решают имена сайтов: мы пишем
\code{wikipedia.org}, а специальный сервер DNS переводит имя в IP-адрес. В домашней сети есть похожий
механизм --- \textbf{mDNS} (англ. \emph{multicast DNS}). Устройство само отвечает на своё имя, которое
всегда оканчивается на \code{.local}. Достаточно одной строки \code{MDNS.begin("esp32-led")}, и
плата станет доступна как \code{esp32-led.local}.

\begin{figure}[H]
\centering
\begin{tikzpicture}[node distance=8mm]
  \node[blk,fill=blkmem,minimum width=24mm] (ph) {Телефон};
  \node[blk,fill=blkcpu,minimum width=24mm] (e) at (9,1.3) {\code{esp32-led}\\ \scriptsize\ttfamily 192.168.1.50};
  \node[blk,fill=blkmem,minimum width=24mm] (n) at (9,-1.3) {Ноутбук};
  \draw[arr,accent] (ph.east) -- node[below,sloped,font=\sffamily\scriptsize]{«Кто esp32-led.local?»} (e.west);
  \draw[arr,accent] (ph.east) -- node[below,sloped,font=\sffamily\scriptsize]{«Кто esp32-led.local?»} (n.west);
  \draw[arr,okgreen] (e.north) |- node[pos=0.75,above,font=\sffamily\scriptsize]{«Это я! Мой адрес 192.168.1.50»} ($(ph.north)+(0,1.6)$) -- (ph.north);
  \node[font=\sffamily\scriptsize,text=gray,right=2mm of n] {молчит};
\end{tikzpicture}
\caption{mDNS: вопрос отправляется всем устройствам сразу, отвечает только то, чьё это имя}
\end{figure}

\begin{tip}
Имена \code{.local} понимают Windows 10/11, macOS, iPhone и Linux. Некоторые старые Android-телефоны
mDNS не поддерживают --- тогда придётся ввести IP-адрес.
\end{tip}


\subsection{Клиент и сервер}

Когда ты открываешь сайт, твой браузер (\textbf{клиент}) отправляет \textbf{запрос} другому компьютеру
(\textbf{серверу}), а тот присылает \textbf{ответ} --- страницу. Это похоже на кафе: посетитель делает
заказ, официант передаёт его на кухню и приносит блюдо. Сегодня кухней будет наша плата: ESP32-S3
станет \textbf{веб-сервером}, а браузер телефона --- клиентом.

\begin{figure}[H]
\centering
\begin{tikzpicture}[node distance=10mm]
  \node[blk,fill=blkmem,minimum width=32mm,minimum height=14mm] (b) {Браузер\\ \scriptsize (клиент)};
  \node[blk,fill=blkcpu,minimum width=32mm,minimum height=14mm,right=55mm of b] (s) {ESP32-S3\\ \scriptsize (веб-сервер)};
  \draw[arr,accent] ($(b.east)+(0,0.3)$) -- node[above,font=\sffamily\scriptsize]{запрос: «дай страницу /»} ($(s.west)+(0,0.3)$);
  \draw[arr,okgreen] ($(s.west)+(0,-0.3)$) -- node[below,font=\sffamily\scriptsize]{ответ: HTML-страница} ($(b.east)+(0,-0.3)$);
\end{tikzpicture}
\caption{Клиент спрашивает --- сервер отвечает. Сам сервер никогда не начинает разговор первым}
\end{figure}

Язык, на котором они разговаривают, называется \textbf{HTTP} (англ. \emph{HyperText Transfer Protocol} ---
протокол передачи гипертекста).

\subsection{Из чего состоит адрес}

\begin{figure}[H]
\centering
\begin{tikzpicture}
  \node[font=\ttfamily\Large,inner sep=0] (u) {%
    \textcolor{gray}{http://}\textcolor{accent}{esp32-led.local}\textcolor{esp}{/level}\textcolor{okgreen!70!black}{?v=30}};
  \draw[decorate,decoration={brace,mirror,amplitude=4pt},gray] ($(u.south west)+(0,-0.1)$) -- ++(1.85,0) node[midway,below=5pt,font=\sffamily\scriptsize,align=center,text=black]{протокол};
  \draw[decorate,decoration={brace,mirror,amplitude=4pt},accent] ($(u.south west)+(1.95,-0.1)$) -- ++(3.9,0) node[midway,below=5pt,font=\sffamily\scriptsize,align=center,text=black]{кому: имя или\\ IP-адрес платы};
  \draw[decorate,decoration={brace,mirror,amplitude=4pt},esp] ($(u.south west)+(5.95,-0.1)$) -- ++(1.5,0) node[midway,below=5pt,xshift=-3mm,font=\sffamily\scriptsize,align=center,text=black]{путь:\\ что хотим};
  \draw[decorate,decoration={brace,mirror,amplitude=4pt},okgreen!70!black] ($(u.south west)+(7.55,-0.1)$) -- ++(1.2,0) node[midway,below=5pt,xshift=4mm,font=\sffamily\scriptsize,align=center,text=black]{параметры:\\ имя = значение};
\end{tikzpicture}
\caption{Адрес (URL) состоит из частей. Параметров может быть несколько: \code{?v=30\&color=red}}
\end{figure}

\subsection{Методы GET и POST}

Каждый запрос начинается со слова-\textbf{метода}, который говорит, чего хочет клиент. Самые важные:

\begin{itemize}
  \item \textbf{GET} («получить») --- «дай мне страницу». Параметры записываются прямо в адресе,
        после знака \code{?}. Так работают обычные ссылки и адресная строка браузера;
  \item \textbf{POST} («отправить») --- «прими от меня данные». Данные передаются в \emph{теле}
        запроса, отдельно от адреса. Так отправляют формы: логин и пароль, сообщение, заказ.
\end{itemize}

\begin{figure}[H]
\centering
\begin{tikzpicture}[
  msg/.style={draw,thick,rounded corners=2pt,align=left,font=\ttfamily\scriptsize,inner sep=5pt,text width=62mm}]
  \node[msg,fill=blkmem!60] (g) {\textbf{GET /level?v=30 HTTP/1.1}\\ Host: esp32-led.local\\ \textcolor{gray}{(тела нет)}};
  \node[msg,fill=blkper!70,right=8mm of g] (p) {\textbf{POST /note HTTP/1.1}\\ Host: esp32.local\\ Content-Type: application/x-www-form-urlencoded\\[2pt]
     \textcolor{esp}{text=Привет,\ плата!}};
  \node[font=\sffamily\small\bfseries,above=1mm of g] {GET --- как открытка};
  \node[font=\sffamily\small\bfseries,above=1mm of p] {POST --- как письмо в конверте};
  \node[font=\sffamily\scriptsize,below=1mm of g,align=center] {данные видны прямо в адресе,\\ их можно сохранить в закладки};
  \node[font=\sffamily\scriptsize,below=1mm of p,align=center] {данные в теле запроса (красным),\\ в адресной строке их не видно};
\end{tikzpicture}
\caption{Так выглядят запросы на самом деле --- это обычный текст}
\end{figure}

\begin{table}[H]
\centering\small
\begin{tabularx}{\textwidth}{@{}l X X@{}}
\toprule
 & \textbf{GET} & \textbf{POST}\\
\midrule
Для чего & получить страницу или данные & отправить данные, что-то изменить\\
Где данные & в адресе: \code{/level?v=30} & в теле запроса\\
Видно в адресной строке & да & нет\\
Можно добавить в закладки & да & нет\\
Повтор при обновлении страницы & безопасен & браузер спросит «отправить повторно?»\\
Объём данных & небольшой (адрес ограничен по длине) & большой (хоть целый файл)\\
\bottomrule
\end{tabularx}
\caption{Сравнение GET и POST}
\end{table}

\begin{tip}
Правило хорошего тона: \textbf{GET} --- только чтобы что-то \emph{узнать}, \textbf{POST} --- чтобы что-то
\emph{изменить}. Иначе поисковый робот или «ускоритель» браузера, заранее открывающий ссылки,
может случайно выключить тебе свет!
\end{tip}

Сервер отвечает \textbf{кодом ответа} --- трёхзначным числом:

\begin{table}[H]
\centering\small
\begin{tabular}{@{}lll@{}}
\toprule
Код & Значение & Когда встречается\\
\midrule
200 & OK --- всё хорошо & страница найдена и отправлена\\
303 & See Other --- смотри другую & «я всё сделал, теперь открой главную»\\
404 & Not Found --- не найдено & адреса нет --- знаменитая «ошибка 404»\\
500 & Internal Server Error & ошибка в программе сервера\\
\bottomrule
\end{tabular}
\caption{Самые частые коды ответа HTTP}
\end{table}


\subsection{Если роутера нет: собственная точка доступа}

ESP32-S3 может сам создать сеть. Для этого вместо подключения к роутеру (\code{WiFi.mode(WIFI\_STA)} и \code{WiFi.begin}) вызывают:

\codecap{Фрагмент: собственная точка доступа}
\codeinput{l8_softap}


\subsection{ESPUI: панель управления без HTML}

Теперь у платы есть сеть, и осталось придумать, как управлять ею с телефона. Писать веб-страницу вручную
сложно (этому посвящён урок~44), поэтому мы воспользуемся библиотекой \textbf{ESPUI}. Вы описываете
панель на C++: «здесь подпись, здесь переключатель, здесь ползунок», --- а библиотека сама создаёт
страницу, отдаёт её браузеру и передаёт события обратно в программу.

\begin{figure}[H]
\centering
\begin{tikzpicture}[node distance=10mm]
  \node[blk,fill=blkmem,minimum width=34mm,minimum height=14mm] (b) {Телефон\\ \scriptsize (браузер)};
  \node[blk,fill=blkcpu,minimum width=34mm,minimum height=14mm,right=60mm of b] (s) {ESP32-S3\\ \scriptsize (ESPUI)};
  \draw[arr,accent] ($(b.east)+(0,0.3)$) -- node[above,font=\sffamily\scriptsize]{событие: «переключатель включён»} ($(s.west)+(0,0.3)$);
  \draw[arr,okgreen] ($(s.west)+(0,-0.3)$) -- node[below,font=\sffamily\scriptsize]{обновление: «подпись изменилась»} ($(b.east)+(0,-0.3)$);
\end{tikzpicture}
\caption{Браузер сообщает о нажатиях, плата присылает новые значения. Страница при этом не перезагружается}
\end{figure}

\paragraph{Установка.} В менеджере библиотек Arduino IDE установите \textsf{ESPUI}, \textsf{ArduinoJson},
\textsf{ESP Async WebServer} и \textsf{Async TCP} (авторы ESP32Async). Проверенные версии и точные
Git-коммиты записаны в \code{examples/web-panels/README.md}; проверку установки и подробности --- в уроке~39.
В PlatformIO всё это подключится само: нужные строки уже есть в \code{platformio.ini} папки \code{sketches/platformio}.

\paragraph{Как подключиться.} Загрузите скетч, откройте Serial Monitor и подключите телефон к сети
\code{S3-Panel} с учебным паролем \code{lesson-05-demo} (для своей платы задайте свой пароль).
Адрес, напечатанный в мониторе (обычно \code{http://192.168.4.1/}), введите в браузере телефона.
Если телефон пишет «без интернета», оставайтесь в этой сети.

\subsubsection{Шаг 1. Первая панель}

Две подписи: одна неизменная, вторая показывает время работы. \code{ESPUI.label(...)} создаёт подпись и
возвращает её номер; по этому номеру мы потом меняем текст функцией \code{updateLabel}.

\codecap{Шаг 1. Панель с двумя подписями}
\codeinput{l05_espui_hello}

\begin{itemize}[leftmargin=*]
  \item \code{WiFi.softAP(имя, пароль)} --- плата сама создаёт сеть (режим из этого урока);
  \item \code{ESPUI.label(заголовок, цвет, текст)} --- элемент-подпись; номер сохраняем, чтобы менять текст;
  \item \code{ESPUI.begin("название")} --- запустить сервер панели; после этого он работает сам, в \code{loop()} его вызывать не нужно.
\end{itemize}

\subsubsection{Шаг 2. Переключатель включает светодиод}

Теперь панель не только показывает, но и \emph{управляет}. Соберите схему из урока~1: GPIO4 $\to$ резистор
\SI{220}{\ohm} $\to$ светодиод $\to$ GND. Переключатель при нажатии вызывает нашу функцию, а она запоминает
желаемое состояние.

\codecap{Шаг 2. Переключатель управляет светодиодом}
\codeinput{l05_espui_led}

Обратите внимание на устройство программы: функция \code{switchCallback} только \emph{записывает желание}
в переменную \code{ledOn}, а выводом управляет \code{loop()}. Так сделано специально: обработчики панели
вызываются из другой задачи FreeRTOS (урок~8), и две задачи не должны одновременно трогать оборудование.
Тип \code{std::atomic<bool>} гарантирует, что запись из одной задачи и чтение из другой не «порвут» значение.

\subsubsection{Шаг 3. Ползунок яркости}

К переключателю добавляем ползунок: он будет задавать яркость, как потенциометр в уроке~4, но на экране
телефона. Подпись ниже показывает, какое число уходит в ШИМ.

\codecap{Шаг 3. Переключатель, ползунок и подпись}
\codeinput{l05_espui_slider}

\begin{itemize}[leftmargin=*]
  \item \code{ESPUI.slider(заголовок, обработчик, цвет, начальное, минимум, максимум)} --- ползунок;
  \item значение приходит в обработчик строкой: \code{sender->value.toInt()} превращает её в число;
  \item \code{sliderContinuous = false} --- отправлять значение, когда палец отпущен, а не сто раз в секунду.
\end{itemize}

\subsubsection{Шаг 4. Своя Wi-Fi-сеть или роутер}

Пока плата сама раздавала сеть. Дома удобнее, когда плата входит в обычную сеть и панель открывается с любого
устройства в ней. В последнем шаге плата сначала пробует подключиться к роутеру, и только если за 15 секунд
не вышло --- создаёт собственную сеть.

\codecap{Шаг 4. Роутер, а при неудаче --- своя сеть}
\codeinput{l05_espui_sta}

Это ровно тот приём, о котором говорилось в разделе «Подключение к сети»: ожидание ограничено по времени, а
запасной вариант гарантирует, что панель всё равно откроется. Подпись внизу показывает уровень сигнала
\code{WiFi.RSSI()} --- отойдите с платой от роутера и посмотрите, как он падает.

\begin{tip}
Чтобы не запоминать IP-адрес, подключите mDNS (раздел выше): \code{MDNS.begin("esp32s3");} --- панель откроется по
адресу \code{http://esp32s3.local/}.
\end{tip}

\subsection{Опыт: гамма-коррекция на ползунке}

Гамма-коррекцию лучше всего «почувствовать» руками: подвигать ползунок и сравнить, как светит
светодиод с поправкой и без неё. Для этого не нужна кнопка и Serial Monitor: плата сама создаёт
панель, а телефон показывает ползунок, переключатель и текущий коэффициент заполнения. Теперь, когда вы умеете собирать панели, добавим к ползунку шага~3 переключатель гаммы из урока~4.

\codecap{ESPUI: яркость светодиода с переключаемой гаммой}
\codeinput{l04_espui_gamma}

Подключите телефон к сети \code{S3-Gamma} (пароль \code{lesson-39-demo}) и откройте адрес из
Serial Monitor. Установите ползунок на 50\,\% и переключайте «Гамма-коррекция»: коэффициент
заполнения упадёт примерно с 127 до 28, а светодиод при этом светит заметно похоже на «половину».
Обратите внимание на устройство программы: callback только записывает число в переменную,
а \code{ledcWrite()} вызывает \code{loop()}.

\subsection{Что дальше с панелями}

Дальше панели будут появляться почти в каждом уроке и усложняться вместе с задачами: ползунок и график к
АЦП (урок~6), вольтметр и радар (уроки~6 и 15), мотор и шаговый двигатель (уроки~16 и 25), цвет ленты
(урок~17). Полный каталог элементов, вкладки и надёжная работа --- в уроке~39. А как устроена страница
«под капотом» --- HTML, CSS и JavaScript --- мы разберём в последнем уроке, когда программировать будет проще.

\begin{trouble}
\begin{itemize}[leftmargin=*]
  \item \textbf{Не подключается к сети} --- сеть 5 ГГц, ошибка в пароле или слабый сигнал (\code{WiFi.RSSI()} ниже $-80$ дБм).
  \item \textbf{Страница не открывается} --- телефон в другой (например, гостевой) сети или адрес набран с ошибкой.
  \item \textbf{Страница открывается очень долго} --- в \code{loop()} есть долгие \code{delay()}, и \code{handleClient()} вызывается редко.
\end{itemize}
\end{trouble}


\begin{summary}
Плата подключается к Wi-Fi как станция или сама раздаёт сеть как точка доступа; панель ESPUI превращает телефон
в пульт: события приходят в обработчики, показания уходят в подписи.

\begin{itemize}[leftmargin=*]
  \item Режимы Wi-Fi: станция, точка доступа, смешанный; только \SI{2.4}{GHz}.
  \item Обработчик панели только записывает желаемое состояние; оборудованием управляет \code{loop()}.
  \item Подключение к роутеру всегда ограничивают по времени и имеют запасной вариант.
\end{itemize}
\end{summary}

\begin{selfcheck}
\begin{enumerate}
  \item Чем режим станции отличается от режима точки доступа?
  \item Почему ожидание подключения к Wi-Fi ограничивают по времени?
  \item Почему обработчик переключателя не должен сам вызывать \code{digitalWrite}?
  \item Чем отличаются GET и POST? Почему POST сам по себе не шифрует данные?
\end{enumerate}
\end{selfcheck}

\begin{tasks}
\begin{enumerate}
\item \lv{1} Загрузите шаг~1, подключитесь к сети платы и откройте панель. Измените текст подписи и цвет элемента.
\item \lv{1} Загрузите шаг~2 и добавьте вторую подпись «Светодиод: включён/выключен», которая обновляется при нажатии (\code{updateLabel}).
\item \lv{2} Добавьте кнопку «Мигнуть»: по нажатию светодиод мигает три раза без \code{delay()} (урок~1, \code{millis()}).
\item \lv{2} На шаге~3 добавьте второй ползунок для максимальной яркости: основной ползунок не должен её превышать.
\item \lv{3} Добавьте подпись с временем работы и числом подключённых к сети телефонов (\code{WiFi.softAPgetStationNum()}).
\item \lv{4} Вместо учебного пароля сделайте так, чтобы плата при первом запуске создавала случайный пароль сети, печатала его в Serial Monitor и показывала в панели только после ввода в поле ввода.
\end{enumerate}
\end{tasks}
